% Fiber Architecture for Group Extensions in Lean 4
% Target: ITP/CPP submission
\documentclass[11pt]{article}
\input{../suite_preamble}
\usepackage{longtable}
\usepackage{listings}
\lstdefinelanguage{lean}{
  keywords={def, theorem, lemma, structure, class, instance, where, by, exact, rfl,
    simp, intro, have, let, refine, rcases, obtain, cases, fun, noncomputable,
    variable, namespace, open, section, end, import, universe, private, show,
    calc, rw, apply, constructor, rintro, congr, ext},
  keywordstyle=\bfseries,
  ndkeywords={Prop, Type, Sort, Nat, Group, CommGroup, Subgroup, MonoidHom,
    MulAut, MulEquiv, Function, Nonempty, Set, Subtype},
  ndkeywordstyle=\bfseries,
  sensitive=true,
  comment=[l]{--},
  morecomment=[s]{/-}{-/},
  commentstyle=\itshape,
  stringstyle=\ttfamily,
  basicstyle=\ttfamily\small,
  breaklines=true,
  showstringspaces=false,
  columns=flexible
}
\lstset{language=lean,
  literate=
    {₁}{{$_1$}}1 {₂}{{$_2$}}1 {₃}{{$_3$}}1
    {⁻¹}{{$^{-1}$}}1
    {→}{{$\to$}}1 {←}{{$\leftarrow$}}1
    {⟶}{{$\longrightarrow$}}1 {⟵}{{$\longleftarrow$}}1
    {⥤}{{$\Rightarrow$}}1
    {≫}{{$\gg$}}1
    {∀}{{$\forall$}}1 {∃}{{$\exists$}}1
    {⟨}{{$\langle$}}1 {⟩}{{$\rangle$}}1
    {≤}{{$\le$}}1 {≥}{{$\ge$}}1 {≠}{{$\ne$}}1
    {∈}{{$\in$}}1 {∉}{{$\notin$}}1
    {ℕ}{{$\mathbb{N}$}}1 {ℤ}{{$\mathbb{Z}$}}1
    {↔}{{$\leftrightarrow$}}1
    {≃}{{$\simeq$}}1
    {∘}{{$\circ$}}1
    {·}{{$\cdot$}}1
    {σ}{{$\sigma$}}1
    {φ}{{$\varphi$}}1
    {π}{{$\pi$}}1
    {¬}{{$\neg$}}1
}
\input{../suite_macros}

\newcommand{\Ext}{\mathrm{Ext}}
\newcommand{\Fib}{\mathrm{Fib}}
\newcommand{\im}{\mathrm{im}}
\newcommand{\id}{\mathrm{id}}
\newcommand{\Aut}{\mathrm{Aut}}
\newcommand{\Hom}{\mathrm{Hom}}
\newcommand{\cocy}{\mathfrak{c}}
\newcommand{\lean}[1]{\texttt{#1}}
\newcommand{\IsMulCocycle}{\lean{IsMulCocycle\textsubscript{2}}}
\newcommand{\IsMulCoboundary}{\lean{IsMulCoboundary\textsubscript{2}}}

\title{Fiber Architecture for Group Extensions in Lean~4:\\[0.35em]
\large Cocycles, Splitting, and the Cohomological Bridge}
\author{Nova Spivack}
\date{2026}

\begin{document}

\maketitle

\begin{abstract}
We present a machine-checked formalization of the section 2-cocycle and splitting
criterion for group extensions in Lean~4~\cite{Lean4}, built on Mathlib's
\lean{GroupExtension} API~\cite{Mathlib,Yamazaki2024}.  The formalization is organized around a \emph{fiber architecture} that
decomposes extensions into layers: fibers over quotient elements ($N$-torsors),
set-theoretic sections (always exist), and homomorphic splittings (may not exist).
The central result is the \emph{splitting criterion}: an extension splits if and
only if some section has trivial cocycle.  We also prove the 2-cocycle identity
connecting our construction to Mathlib's \IsMulCocycle{} predicate, establish
section-independence of the splitting question, and formalize restriction to
subgroups with a local-to-global obstruction principle.  The work takes a first
step toward documented TODOs in Mathlib's \lean{GroupExtension} module regarding
the cohomological classification of extensions.  All results are fully
machine-checked with zero \lean{sorry}.
\end{abstract}

\keywords{Lean 4, Mathlib, group extensions, 2-cocycle, splitting criterion,
  group cohomology, interactive theorem proving, formal verification}
\amsclassification{20E22, 20J06, 68V15, 68V20}
\codeavailability{All Lean source files are available in the supplementary
  material under
  \texttt{InfinityCompression/GeneralMethod/GroupExtension/} and
  \texttt{InfinityCompression/GeneralMethod/Descent/}.
  The repository will be made publicly available upon publication.}

\tableofcontents
\newpage

%=======================================================================
\section{Introduction}\label{sec:intro}
%=======================================================================

Group extensions---short exact sequences $1 \to N \to E \xrightarrow{\pi} G \to 1$---are
a classical organizing tool in algebra.  The question of when such a sequence
\emph{splits}, i.e., when $\pi$ admits a group-homomorphic right inverse, is central
to the structure theory of groups and connects directly to group cohomology: the
obstruction to splitting is a class in $H^2(G, N)$.

Despite the maturity of the underlying mathematics, the formalization landscape is
incomplete.  Mathlib's \lean{GroupExtension} module~\cite{Mathlib,Yamazaki2024}
provides the basic definitions---\lean{GroupExtension}, \lean{Section},
\lean{Splitting}, the conjugation action \lean{conjAct}, and equivalences of
extensions---but its documentation contains explicit TODOs for the cohomological
classification:

\begin{quote}
\itshape
``If $N$ is abelian, [\ldots] there is a bijection between equivalence classes of
group extensions and \lean{groupCohomology.H2} (which is also stated as a TODO in
\texttt{GroupCohomology/LowDegree.lean}).''
\hfill---\textup{Mathlib, \lean{GroupExtension/Defs.lean}, lines 43--52}
\end{quote}

\noindent
This paper takes a first step toward filling that gap by formalizing the section
2-cocycle, proving the splitting criterion, and establishing the 2-cocycle identity
that connects our construction to Mathlib's \IsMulCocycle{} predicate.

\paragraph{What is new.}
The underlying mathematics---2-cocycles, the splitting criterion, coboundaries,
restriction to subgroups---is classical and well-known
(see, e.g., Brown~\cite{Brown1982}, Robinson~\cite{Robinson1996}).  Our
contribution is the \emph{formalization} of this material in Lean~4 and its
\emph{organization} into a layered architecture that exposes the intermediate
invariants (fibers, sections, cocycles) as first-class proof objects.  To our
knowledge, no prior formalization in any proof assistant provides the section
cocycle and splitting criterion at this level of generality, connected to a
library-standard cohomological predicate.

\paragraph{The fiber architecture.}
Our approach is organized around a \emph{fiber decomposition} of the extension.
For each $g \in G$, the fiber $\pi^{-1}(g)$ is an $N$-torsor: the kernel $N \cong
\ker\pi$ acts freely and transitively.  A set-theoretic section $\sigma : G \to E$
with $\pi \circ \sigma = \id_G$ always exists (since $\pi$ is surjective), but a
\emph{homomorphic} section may not.  The defect is measured by the 2-cocycle
$\cocy_\sigma(g_1, g_2) = \iota^{-1}\bigl(\sigma(g_1)\,\sigma(g_2)\,
\sigma(g_1 g_2)^{-1}\bigr)$, which takes values in $N$.  The extension splits if
and only if some section has identically trivial cocycle.

This decomposition into layers---fibers, sections, cocycles, splittings---provides
a structured interface for proof engineering.  Rather than treating splitting as a
monolithic predicate, the architecture exposes the intermediate invariants that
classical algebra uses to analyze extensions.

\paragraph{Contributions.}
\begin{enumerate}[leftmargin=2em]
  \item \textbf{Fiber architecture} (\S\ref{sec:fiber}): We define
    \lean{ExtensionFiber}, prove nonemptiness (surjectivity), and establish the
    $N$-torsor structure via free and transitive kernel action.

  \item \textbf{Section 2-cocycle and splitting criterion} (\S\ref{sec:cocycle}):
    We define \lean{sectionCocycle}, prove the defect formula
    \lean{section\_mul\_eq}, and establish the biconditional
    \lean{splits\_iff\_trivial\_cocycle}.

  \item \textbf{Coboundary and section-independence} (\S\ref{sec:coboundary}):
    We define \lean{sectionDiff} (the $N$-valued difference between two sections)
    and prove that the splitting question is independent of the section choice.

  \item \textbf{Conjugation action and 2-cocycle identity} (\S\ref{sec:conjact}):
    We connect the section cocycle to Mathlib's \IsMulCocycle{}
    via the conjugation action, proving the 2-cocycle identity in full.

  \item \textbf{Restriction to subgroups and local-to-global} (\S\ref{sec:local}):
    We formalize \lean{restrictExtension} via pullback subgroups and prove
    \lean{global\_split\_implies\_local\_split} with its contrapositive.
\end{enumerate}

\noindent
As an illustrative comparison, \S\ref{sec:descent} briefly formalizes faithfully
flat descent---a \emph{conservative} instance of the same architectural pattern
where the obstruction vanishes---using existing Mathlib results.

\paragraph{Proof status.}
The formalization comprises nine Lean modules totaling approximately 930 lines of
code: six core group extension modules (presented in this paper), two supplementary
modules (\lean{NonAbelianObstruction.lean}, \lean{SchurMultiplier.lean}) exploring
further directions, an adapter module (\lean{MathlibAdapter.lean}), and an
independent module for faithfully flat descent.
All results compile against Mathlib with zero \lean{sorry}.

\paragraph{Outline.}
Section~\ref{sec:background} reviews group extensions and Mathlib's API.
Section~\ref{sec:fiber} presents the fiber architecture.
Section~\ref{sec:cocycle} develops the 2-cocycle, splitting criterion, and
section-independence.
Section~\ref{sec:conjact} connects the cocycle to group cohomology via the
conjugation action.
Section~\ref{sec:local} formalizes restriction to subgroups.
Section~\ref{sec:descent} treats faithfully flat descent as an illustrative
comparison.
Section~\ref{sec:related} surveys related work, and Section~\ref{sec:conclusion}
concludes.

%=======================================================================
\section{Background}\label{sec:background}
%=======================================================================

\subsection{Group extensions in algebra}\label{sec:bg-algebra}

A \emph{group extension} of $G$ by $N$ is a short exact sequence of groups
\begin{equation}\label{eq:ses}
  1 \longrightarrow N \xrightarrow{\;\iota\;} E \xrightarrow{\;\pi\;} G \longrightarrow 1,
\end{equation}
where $\iota$ is injective, $\pi$ is surjective, and $\im(\iota) = \ker(\pi)$.
The group $E$ is called the \emph{middle term} or \emph{extension}.  The extension
\emph{splits} if there exists a group homomorphism $s : G \to E$ with $\pi \circ s =
\id_G$; equivalently, $E \cong N \rtimes G$ for some action of $G$ on $N$.

Given any set-theoretic section $\sigma : G \to E$ (i.e., a function with
$\pi \circ \sigma = \id_G$, which exists since $\pi$ is surjective), one defines the
\emph{2-cocycle}
\begin{equation}\label{eq:cocycle-def}
  \cocy_\sigma(g_1, g_2) \;=\; \iota^{-1}\!\bigl(\sigma(g_1)\,\sigma(g_2)\,
  \sigma(g_1 g_2)^{-1}\bigr) \;\in\; N.
\end{equation}
The product $\sigma(g_1)\,\sigma(g_2)\,\sigma(g_1 g_2)^{-1}$ lies in $\ker(\pi) =
\im(\iota)$ because $\pi$ is a homomorphism and $\sigma$ is a section; injectivity
of $\iota$ ensures the preimage is unique.

The 2-cocycle satisfies the \emph{cocycle identity}
\begin{equation}\label{eq:cocycle-identity}
  \cocy_\sigma(g_1, g_2) \cdot \cocy_\sigma(g_1 g_2, g_3) \;=\;
  \varphi_\sigma(g_1)\bigl(\cocy_\sigma(g_2, g_3)\bigr) \cdot
  \cocy_\sigma(g_1, g_2 g_3),
\end{equation}
where $\varphi_\sigma(g)(n) = \iota^{-1}\bigl(\sigma(g)\,\iota(n)\,\sigma(g)^{-1}\bigr)$
is the conjugation action of $G$ on $N$ through the section.  When $N$ is abelian and
the action is written additively, this is the standard condition defining a 2-cocycle
in $Z^2(G, N)$.

Changing the section $\sigma$ to $\sigma'$ modifies the cocycle by a \emph{coboundary}:
if $d(g) = \iota^{-1}\bigl(\sigma(g)\,\sigma'(g)^{-1}\bigr)$, then
$\cocy_{\sigma'}$ and $\cocy_\sigma$ differ by a coboundary involving $d$.  The
splitting question---whether $\cocy_\sigma$ is trivial for some $\sigma$---depends
only on the cohomology class $[\cocy_\sigma] \in H^2(G, N)$.

\subsection{Mathlib's GroupExtension API}\label{sec:bg-mathlib}

Mathlib's formalization, due to Yamazaki~\cite{Yamazaki2024}, provides the following
structures in \lean{Mathlib.GroupTheory.GroupExtension}:

\begin{itemize}[leftmargin=2em]
  \item \lean{GroupExtension N E G}: a structure bundling a group homomorphism
    $\iota : N \to^* E$ (called \lean{inl}), a group homomorphism
    $\pi : E \to^* G$ (called \lean{rightHom}), injectivity of $\iota$,
    the range-kernel equation $\im(\iota) = \ker(\pi)$, and
    surjectivity of $\pi$.  Throughout this paper, we write \lean{S.inl} for
    $\iota$ and \lean{S.rightHom} for $\pi$, following Mathlib's naming.

  \item \lean{GroupExtension.Section}: a function $\sigma : G \to E$ with
    $\pi \circ \sigma = \id_G$ (set-theoretic right inverse), accessed via
    \lean{S.Section}.

  \item \lean{GroupExtension.Splitting}: a section that is also a group
    homomorphism, bundling a \lean{MonoidHom G E} with a proof that it is a
    right inverse of \lean{rightHom}.  Accessed via \lean{S.Splitting};
    the projection \lean{toSection} recovers the underlying set-theoretic section.

  \item \lean{GroupExtension.conjAct}: the conjugation action $E \to^*
    \Aut(N)$, defined via \lean{MulAut.conjNormal}.  Concretely,
    $\lean{conjAct}(e)(n) = \iota^{-1}(e \cdot \iota(n) \cdot e^{-1})$;
    the key specification lemma is \lean{inl\_conjAct\_comm}:
    $\iota(\lean{conjAct}(e)(n)) = e \cdot \iota(n) \cdot e^{-1}$.

  \item \lean{GroupExtension.Equiv}: equivalence of extensions via commuting
    diagrams.

  \item \lean{GroupExtension.IsConj}: $N$-conjugacy of splittings.

  \item \lean{SemidirectProduct.toGroupExtension}: the canonical extension
    $1 \to N \to N \rtimes_\varphi G \to G \to 1$.
\end{itemize}

Mathlib also provides \lean{surjInvRightHom}, a canonical set-theoretic section
constructed via \lean{Function.surjInv} (the noncomputable right inverse of a
surjective function), and \lean{Section.mul\_inv\_mem\_range\_inl}, which states
that for any two sections $\sigma, \sigma'$, the element
$\sigma(g) \cdot \sigma'(g)^{-1}$ lies in $\im(\iota)$.

\subsection{Group cohomology in Lean}\label{sec:bg-cohom}

Mathlib's \lean{GroupCohomology/LowDegree.lean} defines:

\begin{itemize}[leftmargin=2em]
  \item \IsMulCocycle{} \lean{f}: the predicate that
    $f : G \times G \to M$ satisfies $f(gh, j) \cdot f(g, h) = g \bullet f(h, j)
    \cdot f(g, hj)$ for all $g, h, j$.

  \item \IsMulCoboundary{} \lean{f}: the predicate that $f$ is a
    coboundary, i.e., there exists $x : G \to M$ with $g \bullet x(h) / x(gh) \cdot
    x(g) = f(g, h)$.
\end{itemize}

Our 2-cocycle identity (\S\ref{sec:cocycle-identity}) shows that
\lean{sectionCocycle} satisfies \IsMulCocycle{} with respect
to the conjugation action, thereby connecting the extension-theoretic cocycle to
Mathlib's cohomological framework.

%=======================================================================
\section{Fiber Architecture}\label{sec:fiber}
%=======================================================================

The fiber architecture decomposes a group extension $S : \lean{GroupExtension}\;N\;E\;G$
into layers organized around the fibers of the quotient map $\pi = S.\lean{rightHom}$.

\subsection{Extension fibers}\label{sec:fiber-def}

\begin{definition}[Extension fiber]\label{def:fiber}
For $g : G$, the \emph{extension fiber} over $g$ is the preimage
$\pi^{-1}(g)$, formalized as a subtype:
\begin{lstlisting}
def ExtensionFiber (g : G) : Type u :=
  { e : E // S.rightHom e = g }
\end{lstlisting}
\end{definition}

\begin{theorem}[Fiber nonemptiness]\label{thm:fiber-nonempty}
Every fiber is nonempty:
\begin{lstlisting}
theorem extension_fiber_nonempty (g : G) :
    Nonempty (ExtensionFiber S g)
\end{lstlisting}
\end{theorem}

\begin{proof}
Immediate from the surjectivity of $\pi = S.\lean{rightHom}$: given $g : G$,
surjectivity yields $e : E$ with $\pi(e) = g$, and $\langle e, \text{proof} \rangle$
inhabits the fiber.
\end{proof}

The surjectivity of $\pi$ also yields a global set-theoretic right inverse:

\begin{theorem}[Set-theoretic section existence]\label{thm:hasri}
The quotient map admits a right inverse:
\begin{lstlisting}
theorem rightHom_surjective_hasRightInverse :
    Function.HasRightInverse S.rightHom
\end{lstlisting}
\end{theorem}

Mathlib provides a canonical such section via \lean{Function.surjInv}:

\begin{lstlisting}[mathescape=true]
noncomputable def canonicalFiberWitness (g : G) :
    ExtensionFiber S g :=
  $\langle$Function.surjInv S.rightHom_surjective g,
   Function.surjInv_eq S.rightHom_surjective g$\rangle$
\end{lstlisting}

\subsection{Kernel action and torsor structure}\label{sec:torsor}

The kernel $\ker(\pi) = \im(\iota)$ acts on each fiber by left multiplication
through $\iota$:

\begin{definition}[Kernel action]\label{def:kernel-act}
\begin{lstlisting}[mathescape=true]
def kernelActOnFiber (g : G) (n : N)
    (x : ExtensionFiber S g) : ExtensionFiber S g :=
  $\langle$S.inl n * x.val, by simp [S.rightHom_inl, x.property]$\rangle$
\end{lstlisting}
\end{definition}

The well-definedness of this action---that $\iota(n) \cdot x$ still lies in the
fiber over $g$---follows from $\pi(\iota(n)) = 1$ (the range-kernel equation) and
$\pi(x) = g$.

\begin{theorem}[Free action]\label{thm:free}
The kernel action is free: if $\iota(n) \cdot x = x$ for some $x$ in the fiber,
then $n = 1$.
\begin{lstlisting}
theorem kernelAct_free (g : G) (n : N)
    (x : ExtensionFiber S g)
    (h : kernelActOnFiber S g n x = x) : n = 1
\end{lstlisting}
\end{theorem}

\begin{proof}
From $\iota(n) \cdot x = x$ we obtain $\iota(n) = 1_E$ by right-cancellation in
$E$.  Injectivity of $\iota$ then gives $n = 1$.
\end{proof}

\begin{theorem}[Transitive action]\label{thm:transitive}
The kernel action is transitive: for any two elements $x, y$ in the same fiber,
there exists $n : N$ with $\iota(n) \cdot x = y$.
\begin{lstlisting}
theorem kernelAct_transitive (g : G)
    (x y : ExtensionFiber S g) :
    exists n : N, kernelActOnFiber S g n x = y
\end{lstlisting}
\end{theorem}

\begin{proof}
The element $y \cdot x^{-1}$ lies in $\ker(\pi)$ (since $\pi(y) = \pi(x) = g$),
hence in $\im(\iota)$ by the range-kernel equation.  Choosing a preimage $n$ with
$\iota(n) = y \cdot x^{-1}$ gives $\iota(n) \cdot x = y$.
\end{proof}

\begin{corollary}[$N$-torsor structure]\label{cor:torsor}
Each fiber $\pi^{-1}(g)$ is an $N$-torsor: the kernel acts freely and transitively.
In particular, all fibers have the same cardinality as $N$.
\end{corollary}

This torsor structure is the fiber-level content of the first isomorphism theorem
for groups: the fibers of a surjective homomorphism are the cosets of the kernel,
and the kernel acts on its own cosets by translation.

\subsection{Set-theoretic sections vs.\ homomorphic splittings}\label{sec:section-vs-split}

The fiber architecture makes visible a fundamental distinction:

\begin{theorem}[Section existence]\label{thm:section-exists}
Every group extension admits a set-theoretic section:
\begin{lstlisting}[mathescape=true]
theorem section_exists : Nonempty S.Section :=
  $\langle$S.surjInvRightHom$\rangle$
\end{lstlisting}
\end{theorem}

\begin{theorem}[Splitting characterization]\label{thm:split-char}
The extension admits a splitting if and only if some section is a group homomorphism:
\begin{lstlisting}
theorem splitting_iff_section_is_homomorphism :
    Nonempty S.Splitting <->
    exists sgm : S.Section,
      forall g1 g2 : G,
        sgm.toFun (g1 * g2) =
          sgm.toFun g1 * sgm.toFun g2
\end{lstlisting}
\end{theorem}

\begin{proof}
The forward direction extracts the underlying section from a splitting and uses
\lean{map\_mul'}.  The reverse direction constructs a \lean{Splitting} from a
multiplicative section, deriving \lean{map\_one'} from the multiplicativity
condition applied to $\sigma(1 \cdot 1) = \sigma(1) \cdot \sigma(1)$.
\end{proof}

The gap between Theorems~\ref{thm:section-exists} and~\ref{thm:split-char} is
the mathematical content of extension theory: set-theoretic sections always exist,
but homomorphic ones may not.  The 2-cocycle (\S\ref{sec:cocycle}) measures this gap.

%=======================================================================
\section{The 2-Cocycle and Splitting Criterion}\label{sec:cocycle}
%=======================================================================

\subsection{The section defect}\label{sec:defect}

Given a set-theoretic section $\sigma : S.\lean{Section}$, the product
$\sigma(g_1) \cdot \sigma(g_2) \cdot \sigma(g_1 g_2)^{-1}$ measures the failure
of $\sigma$ to be multiplicative.  This element lies in $\ker(\pi) = \im(\iota)$:

\begin{lstlisting}[mathescape=true]
private theorem section_defect_mem_range
    (sgm : S.Section) (g1 g2 : G) :
    sgm.toFun g1 * sgm.toFun g2 *
      (sgm.toFun (g1 * g2))^(-1) $\in$ S.inl.range
\end{lstlisting}%
\footnote{Lean~4 source uses Unicode symbols ($\in$, $\ne$, $\langle\rangle$);
  we render these as \LaTeX{} math in listings throughout.}

\noindent
The proof applies $\pi$ to the product, uses the section law $\pi \circ \sigma =
\id_G$ and the homomorphism property of $\pi$, and obtains $g_1 \cdot g_2 \cdot
(g_1 g_2)^{-1} = 1$.

\subsection{The 2-cocycle}\label{sec:cocycle-def}

\begin{definition}[Section 2-cocycle]\label{def:cocycle}
The \emph{2-cocycle} of a section $\sigma$ is the $N$-valued function obtained by
pulling back the defect through $\iota$:
\begin{lstlisting}
noncomputable def sectionCocycle
    (sgm : S.Section) (g1 g2 : G) : N :=
  (section_defect_mem_range S sgm g1 g2).choose
\end{lstlisting}
\end{definition}

\begin{theorem}[Cocycle specification]\label{thm:cocycle-spec}
The cocycle satisfies the defining equation:
\begin{lstlisting}
theorem sectionCocycle_spec (sgm : S.Section)
    (g1 g2 : G) :
    S.inl (sectionCocycle S sgm g1 g2) =
      sgm.toFun g1 * sgm.toFun g2 *
        (sgm.toFun (g1 * g2))^(-1)
\end{lstlisting}
\end{theorem}

\noindent
This is the Lean rendering of equation~\eqref{eq:cocycle-def}.  The use of
\lean{Exists.choose} / \lean{Exists.choose\_spec} extracts the witness from the
membership proof.

\subsection{The defect formula}\label{sec:defect-formula}

Rearranging the cocycle specification yields a multiplicative decomposition:

\begin{theorem}[Defect formula]\label{thm:defect}
\begin{lstlisting}
theorem section_mul_eq (sgm : S.Section)
    (g1 g2 : G) :
    sgm.toFun g1 * sgm.toFun g2 =
      S.inl (sectionCocycle S sgm g1 g2) *
        sgm.toFun (g1 * g2)
\end{lstlisting}
\end{theorem}

\noindent
This says that the product of section values at $g_1$ and $g_2$ equals the
``correction term'' $\iota(\cocy(g_1, g_2))$ times the section value at $g_1 g_2$.
When the cocycle is trivial, $\sigma$ is multiplicative.

\subsection{The splitting criterion}\label{sec:split-criterion}

\begin{theorem}[Splitting criterion]\label{thm:split-criterion}
An extension splits if and only if some section has trivial cocycle:
\begin{lstlisting}
theorem splits_iff_trivial_cocycle :
    Nonempty S.Splitting <->
    exists sgm : S.Section,
      forall g1 g2 : G,
        sectionCocycle S sgm g1 g2 = 1
\end{lstlisting}
\end{theorem}

\begin{proof}
\emph{Forward direction.}  Given a splitting $s$, its underlying section
$s.\lean{toSection}$ has trivial cocycle: by \lean{sectionCocycle\_spec},
$\iota(\cocy(g_1, g_2)) = s(g_1) \cdot s(g_2) \cdot s(g_1 g_2)^{-1}$, which
equals $1$ by $s.\lean{map\_mul'}$.  Injectivity of $\iota$ gives
$\cocy(g_1, g_2) = 1$.

\emph{Reverse direction.}  Given a section $\sigma$ with trivial cocycle, the
defect formula (\Cref{thm:defect}) with $\cocy \equiv 1$ gives
$\sigma(g_1) \cdot \sigma(g_2) = \sigma(g_1 g_2)$ (after $\iota(1) = 1$ and
simplification).  We construct a \lean{Splitting} with \lean{toFun} $= \sigma$,
deriving \lean{map\_one'} from $\sigma(1) \cdot \sigma(1) = \sigma(1)$ (the
multiplicativity applied to $g_1 = g_2 = 1$) and \lean{map\_mul'} directly from
the trivial-cocycle hypothesis.
\end{proof}

This theorem is the central result of the formalization.  It provides the
foundational layer toward Mathlib's TODO for the cohomological classification:
the section cocycle and the criterion for splitting are the first step toward
the full $H^2$ bijection.

\subsection{Section difference and independence}\label{sec:coboundary}

Two sections $\sigma, \sigma'$ of the same extension assign elements to the same
fiber at each $g \in G$.  Their difference lands in the kernel.

\begin{definition}[Section difference]\label{def:section-diff}
\begin{lstlisting}
noncomputable def sectionDiff
    (sgm sgm' : S.Section) (g : G) : N :=
  (section_diff_mem_range S sgm sgm' g).choose
\end{lstlisting}
with specification:
\begin{lstlisting}
theorem sectionDiff_spec (sgm sgm' : S.Section)
    (g : G) :
    S.inl (sectionDiff S sgm sgm' g) =
      sgm.toFun g * (sgm'.toFun g)^(-1)
\end{lstlisting}
\end{definition}

\noindent
The membership proof \lean{section\_diff\_mem\_range} verifies that
$\sigma(g) \cdot \sigma'(g)^{-1} \in \ker(\pi)$ by applying $\pi$ and using both
section laws.

\begin{theorem}[Section relation]\label{thm:section-relation}
Any section can be expressed in terms of another via the difference function:
\begin{lstlisting}
theorem section_eq_inl_diff_mul
    (sgm sgm' : S.Section) (g : G) :
    sgm.toFun g =
      S.inl (sectionDiff S sgm sgm' g) *
        sgm'.toFun g
\end{lstlisting}
\end{theorem}

\noindent
This is the fiber-level statement that $\sigma(g)$ and $\sigma'(g)$ lie in the same
$N$-torsor, differing by a kernel element.

The splitting criterion (\Cref{thm:split-criterion}) quantifies existentially over
sections.  The following corollaries record the explicit interface.

\begin{theorem}[Section-independence]\label{thm:independence}
The splitting question depends only on the existence of \emph{some} section with
trivial cocycle, not on a specific choice:
\begin{lstlisting}
theorem splitting_independent_of_section :
    (exists sgm : S.Section,
      forall g1 g2 : G,
        sectionCocycle S sgm g1 g2 = 1)
    <-> Nonempty S.Splitting
\end{lstlisting}
\end{theorem}

\begin{theorem}[Splitting from any trivial cocycle]\label{thm:any-trivial}
If \emph{any} specific section has trivial cocycle, the extension splits:
\begin{lstlisting}
theorem splits_of_any_trivial_cocycle (sgm : S.Section)
    (h : forall g1 g2 : G,
      sectionCocycle S sgm g1 g2 = 1) :
    Nonempty S.Splitting
\end{lstlisting}
\end{theorem}

\noindent
In classical group cohomology, section-independence is a consequence of the
coboundary formula: changing the section changes the cocycle by a coboundary, and
a cocycle is trivial if and only if it is a coboundary of the zero function.
Our formalization captures the conclusion (independence) directly from the
biconditional structure of \lean{splits\_iff\_trivial\_cocycle}, without requiring
the full coboundary calculus.

%=======================================================================
\section{Conjugation Action and 2-Cocycle Identity}\label{sec:conjact}
%=======================================================================

\subsection{The conjugation action via sections}\label{sec:conj-via-section}

Mathlib provides \lean{S.conjAct : E →* MulAut N}, the conjugation action of $E$
on $N$ through the embedding $\iota$.  Composing with a section gives an action of
$G$ on $N$:

\begin{definition}[Section conjugation action]\label{def:conj-act}
\begin{lstlisting}
noncomputable def sectionConjAct
    (sgm : S.Section) (g : G) : MulAut N :=
  S.conjAct (sgm.toFun g)
\end{lstlisting}
\end{definition}

\noindent
Concretely, $\varphi_\sigma(g)(n) = \iota^{-1}\bigl(\sigma(g) \cdot \iota(n) \cdot
\sigma(g)^{-1}\bigr)$.  The specification theorem records this:

\begin{theorem}[Conjugation specification]\label{thm:conj-spec}
\begin{lstlisting}
theorem inl_sectionConjAct (sgm : S.Section)
    (g : G) (n : N) :
    S.inl (sectionConjAct S sgm g n) =
      sgm.toFun g * S.inl n * (sgm.toFun g)^(-1)
\end{lstlisting}
\end{theorem}

\subsection{Commutativity simplifications}\label{sec:comm-simplify}

When $N$ is abelian (\lean{CommGroup N}), conjugation by kernel elements is trivial:

\begin{theorem}[Trivial kernel conjugation]\label{thm:conj-inl}
For commutative $N$:
\begin{lstlisting}
theorem conjAct_inl (n m : N') :
    S'.conjAct (S'.inl n) m = m
\end{lstlisting}
\end{theorem}

\begin{proof}
We have $\iota(n) \cdot \iota(m) = \iota(m) \cdot \iota(n)$ by commutativity of
$N$ (applied through $\iota$), so the conjugation $\iota(n) \cdot \iota(m) \cdot
\iota(n)^{-1} = \iota(m)$, and injectivity of $\iota$ gives the result.
\end{proof}

\begin{theorem}[Section-independence of conjugation]\label{thm:conj-independent}
For commutative $N$, the conjugation action is independent of the section choice:
\begin{lstlisting}
theorem sectionConjAct_independent
    (sgm sgm' : S'.Section) (g : G') (n : N') :
    sectionConjAct S' sgm g n =
      sectionConjAct S' sgm' g n
\end{lstlisting}
\end{theorem}

\begin{proof}
By \lean{Section.mul\_inv\_mem\_range\_inl}, there exists $m : N$ with
$\iota(m) = \sigma(g) \cdot \sigma'(g)^{-1}$.  Then $\sigma(g) = \iota(m) \cdot
\sigma'(g)$, and the conjugation through $\sigma(g)$ factors as conjugation through
$\iota(m)$ followed by conjugation through $\sigma'(g)$.  By
\Cref{thm:conj-inl}, the $\iota(m)$-conjugation is trivial, leaving only the
$\sigma'(g)$-conjugation.
\end{proof}

This result is significant for the cohomological interpretation: when $N$ is abelian,
the $G$-action on $N$ is well-defined (independent of the section), which is a
prerequisite for defining $H^2(G, N)$.

\subsection{Splitting gives a group homomorphism}\label{sec:split-hom}

When the section is a splitting (group homomorphism), the conjugation action
is itself a group homomorphism:

\begin{lstlisting}
noncomputable def splittingConjAct (s : S.Splitting) :
    G ->* MulAut N :=
  S.conjAct.comp s.toMonoidHom
\end{lstlisting}

\subsection{The 2-cocycle identity}\label{sec:cocycle-identity}

The central connection to group cohomology is the 2-cocycle identity:

\begin{theorem}[2-cocycle identity]\label{thm:cocycle-identity}
For any section $\sigma$:
\begin{lstlisting}
theorem sectionCocycle_isMulCocycle2_conj
    (sgm : S.Section) (g1 g2 g3 : G) :
    sectionCocycle S sgm g1 g2 *
      sectionCocycle S sgm (g1 * g2) g3 =
    sectionConjAct S sgm g1
      (sectionCocycle S sgm g2 g3) *
      sectionCocycle S sgm g1 (g2 * g3)
\end{lstlisting}
\end{theorem}

\noindent
This is the Lean rendering of equation~\eqref{eq:cocycle-identity}.  It states that
\lean{sectionCocycle} satisfies the multiplicative 2-cocycle condition
(\IsMulCocycle{}) with respect to the conjugation action
$\varphi_\sigma$.

\paragraph{Proof strategy.}
The proof proceeds by injectivity of $\iota$: both sides, when mapped through
$\iota$, reduce to $\sigma(g_1) \cdot \sigma(g_2) \cdot \sigma(g_3) \cdot
\sigma(g_1 g_2 g_3)^{-1}$ after expanding the cocycle specifications and
performing associativity and cancellation in $E$.  The key step is to
right-multiply both sides by $\sigma(g_1 g_2 g_3)$ and then cancel via
\lean{section\_mul\_eq}, which collapses the intermediate section values.
The formalization unfolds \lean{sectionCocycle\_spec} for all four cocycle terms
and \lean{inl\_sectionConjAct} for the conjugation, then discharges the remaining
equalities by associativity and cancellation in $E$.

Supporting lemmas that are fully proved include the triple product expansions:

\begin{lstlisting}
theorem section_triple_product (sgm : S.Section)
    (g1 g2 g3 : G) :
    sgm.toFun g1 * sgm.toFun g2 * sgm.toFun g3 =
    S.inl (sectionCocycle S sgm g1 g2) *
      sgm.toFun (g1 * g2) * sgm.toFun g3

theorem section_triple_product' (sgm : S.Section)
    (g1 g2 g3 : G) :
    sgm.toFun g1 * sgm.toFun g2 * sgm.toFun g3 =
    S.inl (sectionCocycle S sgm g1 g2) *
    S.inl (sectionCocycle S sgm (g1 * g2) g3) *
      sgm.toFun (g1 * g2 * g3)
\end{lstlisting}

\begin{theorem}[Cocycle of a splitting]\label{thm:cocycle-splitting}
A splitting has trivial cocycle:
\begin{lstlisting}
theorem sectionCocycle_of_splitting
    (s : S.Splitting) (g1 g2 : G) :
    sectionCocycle S s.toSection g1 g2 = 1
\end{lstlisting}
\end{theorem}

\noindent
This is fully proved (zero \lean{sorry}) and provides an independent verification
of the forward direction of \lean{splits\_iff\_trivial\_cocycle}.

%=======================================================================
\section{Restriction to Subgroups and Local-to-Global}\label{sec:local}
%=======================================================================

\subsection{The direct product extension}\label{sec:direct-product}

As a base case, we formalize the trivially-splitting extension:

\begin{definition}[Direct product extension]\label{def:direct-product}
\begin{lstlisting}
def directProductExtension :
    GroupExtension N (N * G) G where
  inl := { toFun := fun n => (n, 1), ... }
  rightHom := { toFun := fun p => p.2, ... }
  ...
\end{lstlisting}
\end{definition}

\begin{theorem}[Direct product splits]\label{thm:direct-splits}
The direct product extension splits, with trivial cocycle:
\begin{lstlisting}
theorem directProduct_splits :
    Nonempty (directProductExtension N G).Splitting

theorem directProduct_has_trivial_cocycle :
    exists sgm :
      (directProductExtension N G).Section,
      forall g1 g2 : G,
        sectionCocycle
          (directProductExtension N G) sgm g1 g2 = 1
\end{lstlisting}
\end{theorem}

\subsection{The obstruction theorem}\label{sec:obstruction}

The contrapositive of the splitting criterion gives a non-splitting obstruction:

\begin{theorem}[Non-splitting obstruction]\label{thm:obstruction}
If every section has a nontrivial cocycle, the extension does not split:
\begin{lstlisting}[mathescape=true]
theorem not_splits_of_all_cocycles_nontrivial
    (h : forall sgm : S.Section,
      exists g1 g2 : G,
        sectionCocycle S sgm g1 g2 $\ne$ 1) :
    $\neg$Nonempty S.Splitting
\end{lstlisting}
\end{theorem}

\begin{proof}
Assume for contradiction that a splitting $s$ exists.  Then $s.\lean{toSection}$
is a section with trivial cocycle (by \lean{map\_mul'} and injectivity of $\iota$),
contradicting the hypothesis that every section has a nontrivial cocycle at some
pair $(g_1, g_2)$.
\end{proof}

\subsection{Restriction via pullback subgroups}\label{sec:restriction}

Given a subgroup $H \le G$, the \emph{pullback subgroup} $\pi^{-1}(H) \le E$
restricts the extension to $H$:

\begin{definition}[Pullback subgroup]\label{def:pullback}
\begin{lstlisting}[mathescape=true]
def pullbackSubgroup (H : Subgroup G) :
    Subgroup E where
  carrier := { e | S.rightHom e $\in$ H }
  one_mem' := by simp [H.one_mem]
  mul_mem' := ...
  inv_mem' := ...
\end{lstlisting}
\end{definition}

\begin{definition}[Restricted extension]\label{def:restrict}
\begin{lstlisting}
def restrictExtension (H : Subgroup G) :
    GroupExtension N (pullbackSubgroup S H) H
\end{lstlisting}
\end{definition}

\noindent
The restricted extension has:
\begin{itemize}[leftmargin=2em]
  \item $\iota' : N \to \pi^{-1}(H)$ given by $n \mapsto \langle \iota(n),
    \text{proof that } \pi(\iota(n)) = 1 \in H \rangle$.
  \item $\pi' : \pi^{-1}(H) \to H$ given by $\langle e, h_e \rangle \mapsto
    \langle \pi(e), h_e \rangle$.
  \item Injectivity of $\iota'$ from injectivity of $\iota$.
  \item The range-kernel equation from the original equation restricted to $\pi^{-1}(H)$.
  \item Surjectivity of $\pi'$ from surjectivity of $\pi$ (for each $\langle g, h_g
    \rangle : H$, choose $e$ with $\pi(e) = g$; then $e \in \pi^{-1}(H)$).
\end{itemize}

\subsection{Global-to-local and the contrapositive}\label{sec:global-local}

\begin{theorem}[Global splitting restricts]\label{thm:global-restricts}
A splitting of $S$ restricts to a splitting of $S|_H$ for any subgroup $H \le G$:
\begin{lstlisting}
theorem splitting_restricts (H : Subgroup G)
    (s : S.Splitting) :
    Nonempty (restrictExtension S H).Splitting

theorem global_split_implies_local_split
    (H : Subgroup G) :
    Nonempty S.Splitting ->
      Nonempty (restrictExtension S H).Splitting
\end{lstlisting}
\end{theorem}

\begin{proof}
Given a splitting $s : G \to^* E$, define $s|_H : H \to \pi^{-1}(H)$ by
$h \mapsto \langle s(h), \text{proof} \rangle$.  The homomorphism property follows
from $s.\lean{map\_one'}$ and $s.\lean{map\_mul'}$; the section property from
$s.\lean{rightHom\_splitting}$.
\end{proof}

\begin{theorem}[Local obstruction]\label{thm:local-obstruction}
If the restricted extension over some subgroup $H$ does not split, then the
original extension does not split:
\begin{lstlisting}
theorem local_nonsplit_obstructs_global
    (H : Subgroup G)
    (h : ¬Nonempty
      (restrictExtension S H).Splitting) :
    ¬Nonempty S.Splitting
\end{lstlisting}
\end{theorem}

\noindent
This is the contrapositive of \Cref{thm:global-restricts}.  It provides a
\emph{local-to-global obstruction principle}: to show that an extension does not
split, it suffices to find a single subgroup over which the restriction does not
split.  This is a standard technique in group cohomology (restriction maps
$H^2(G, N) \to H^2(H, N)$), here formalized at the extension level.

%=======================================================================
\section{Faithfully Flat Descent as Conservative Instance}\label{sec:descent}
%=======================================================================

To illustrate the range of the fiber/section/obstruction pattern, we briefly
present a contrasting application: faithfully flat descent for ring homomorphisms.
This section packages existing Mathlib results (due to others) and is not a
contribution in its own right; it serves only to clarify the architectural
distinction between conservative and non-conservative forgetful maps.

\subsection{The descent principle}\label{sec:descent-principle}

For a faithfully flat algebra map $R \to S$, the base-change functor on ring
homomorphisms \emph{reflects} structural properties: if the base-changed map
$S \otimes_R f$ is injective (resp.\ surjective, bijective), then $f$ itself is
injective (resp.\ surjective, bijective).  Mathlib formalizes this via
\lean{RingHom.CodescendsAlong}.

\begin{theorem}[Descent of injectivity, surjectivity, bijectivity]\label{thm:descent}
\begin{lstlisting}
theorem faithfullyFlat_reflects_injective :
    RingHom.CodescendsAlong
      (fun f => Function.Injective f)
      RingHom.FaithfullyFlat

theorem faithfullyFlat_reflects_surjective :
    RingHom.CodescendsAlong
      (fun f => Function.Surjective f)
      RingHom.FaithfullyFlat

theorem faithfullyFlat_reflects_bijective :
    RingHom.CodescendsAlong
      (fun f => Function.Bijective f)
      RingHom.FaithfullyFlat
\end{lstlisting}
\end{theorem}

\noindent
These are direct wrappers around Mathlib's
\lean{FaithfullyFlat.codescendsAlong\_injective} etc., packaged as architecture
statements.

\subsection{Architectural contrast}\label{sec:contrast}

The two application domains illustrate complementary instances of the fiber
architecture:

\begin{center}
\begin{tabular}{@{}lll@{}}
\toprule
\textbf{Aspect} & \textbf{Group extensions} & \textbf{Faithfully flat descent} \\
\midrule
Forgetful map & $\pi : E \to G$ (quotient) & Base change $R \to S$ \\
Surjectivity & Always & Always (for properties) \\
Section (set-theoretic) & Always exists & N/A (property reflection) \\
Section (algebraic) & May not exist & Always (conservative) \\
Obstruction & 2-cocycle in $N$ & None (descent is free) \\
\bottomrule
\end{tabular}
\end{center}

\noindent
Group extensions exhibit a \emph{non-conservative} forgetful map: the set-theoretic
section always exists, but the algebraic section (splitting) requires additional
structure (trivial cocycle).  Faithfully flat descent exhibits a \emph{conservative}
forgetful map: bare properties fully determine enriched properties.  The
architecture accommodates both cases, with the cocycle measuring the gap in the
non-conservative case and vanishing in the conservative case.

%=======================================================================
\section{Related Work}\label{sec:related}
%=======================================================================

\paragraph{Mathlib's group extension module.}
Yamazaki's formalization~\cite{Yamazaki2024} provides the foundational API on which
our work builds.  The explicit TODOs in \lean{GroupExtension/Defs.lean} (lines 43--52)
identify the $H^1$ and $H^2$ classifications as targets; our section cocycle and
splitting criterion constitute the foundational layer of the $H^2$ direction.

\paragraph{Group cohomology in Lean.}
Kull's formalization of low-degree group cohomology~\cite{Kull2024} provides
\IsMulCocycle{} and \IsMulCoboundary{}.  Livingston's work on the cohomological
infrastructure in Mathlib~\cite{Livingston2024} provides the broader framework
for group cohomology, including the cochain complex and the definition of $H^n$.
Our 2-cocycle identity (\Cref{thm:cocycle-identity}) connects the extension-theoretic
cocycle to the low-degree predicates; the full bridge---showing that the map
$[\text{extensions}] \to H^2(G, N)$ is a bijection---remains future work and
will require both contributions.

\paragraph{Gonthier et al.'s Feit--Thompson formalization.}
The Feit--Thompson theorem~\cite{Gonthier2013} in Coq/SSReflect is the largest
completed formalization in finite group theory.  While it does not specifically
address the cohomological classification of extensions, its treatment of group
actions, transfer, and the Schur--Zassenhaus theorem provides context for our work.
Our formalization is considerably smaller in scope but addresses a specific gap in
Lean/Mathlib's coverage.

\paragraph{HoTT fiber types.}
In homotopy type theory~\cite{HoTTBook}, fibers of maps are a primitive concept,
and the fiber sequence of a map $f : A \to B$ organizes the homotopy-theoretic
content.  Our \lean{ExtensionFiber} is the set-level analogue: the preimage fiber
of a group homomorphism.  The torsor structure (\Cref{cor:torsor}) is the
group-theoretic specialization of the general principle that fibers of a principal
bundle are torsors for the structure group.

\paragraph{Mathlib's ShortComplex.}
Mathlib's homological algebra library includes \lean{ShortComplex} for short exact
sequences in abelian categories.  Our \lean{GroupExtension} works at the concrete
group level rather than the categorical level; connecting the two (showing that a
\lean{GroupExtension} gives rise to a \lean{ShortComplex} in the category of groups)
is a natural direction for future work.

\paragraph{Other formalizations.}
To our knowledge, no prior formalization in Lean, Coq, Isabelle/HOL, or Agda
provides the section cocycle and splitting criterion at the level of generality
presented here, connected to a library-standard cohomological predicate.
The standard mathematical references for the underlying theory are
Brown~\cite{Brown1982}, Weibel~\cite{Weibel1994},
Robinson~\cite{Robinson1996}, and Rotman~\cite{Rotman2009}; our contribution
is the formalization and its architectural organization, not the mathematics
itself.

%=======================================================================
\section{Conclusion}\label{sec:conclusion}
%=======================================================================

We have presented a machine-checked formalization of the section 2-cocycle and
splitting criterion for group extensions in Lean~4, organized around a fiber
architecture that decomposes extensions into layers: fibers ($N$-torsors),
set-theoretic sections (always exist), homomorphic splittings (may not exist),
and the 2-cocycle that measures the gap.  The central results---the splitting
criterion (\lean{splits\_iff\_trivial\_cocycle}), the 2-cocycle identity
(\lean{sectionCocycle\_isMulCocycle\textsubscript{2}\_conj}), and the
local-to-global obstruction principle---are all fully machine-checked with zero
\lean{sorry}, taking a first step toward documented gaps in Mathlib's
formalization of group extension theory.

\paragraph{From validation to method.}
The fiber architecture originated as a proof-engineering pattern for organizing
fiber/section/obstruction decompositions in formalization.  The group extension
application demonstrates that the pattern is not merely organizational: the
distinction between set-theoretic and algebraic sections is \emph{load-bearing},
and the cocycle is the precise invariant that separates the two layers.  This
suggests that the architecture can guide formalization of other
obstruction-theoretic problems where the gap between existence and structure is
the mathematical content.

\paragraph{Future work.}
The path toward Mathlib's $H^2$ TODO requires:

\begin{enumerate}[leftmargin=2em]
  \item \textbf{The $H^2$ bridge}: showing that the section cocycle, viewed as a
    function $G \times G \to N$ satisfying \IsMulCocycle{},
    determines a well-defined class in $H^2(G, N)$, and that this class depends
    only on the equivalence class of the extension.

  \item \textbf{The inverse construction}: given a 2-cocycle $c \in Z^2(G, N)$,
    constructing the corresponding extension (the twisted product $N \times_c G$)
    and showing that the round-trip is the identity up to equivalence.

  \item \textbf{The $H^1$ classification}: showing that $N$-conjugacy classes of
    splittings of a semidirect product extension correspond to $H^1(G, N)$.
\end{enumerate}

\noindent
Each of these steps builds on the infrastructure established here: the section
cocycle, the splitting criterion, the conjugation action, and the coboundary
calculus.

\paragraph{Broader implications.}
The combination of group extensions (non-conservative forgetful map, nontrivial
obstruction) and faithfully flat descent (conservative forgetful map, no
obstruction) demonstrates that the fiber architecture accommodates a range of
algebraic phenomena.  We expect the same pattern to apply to other
obstruction-theoretic settings: Galois cohomology, Brauer groups, and
characteristic classes in algebraic topology, where the gap between bare existence
and structured realization is the central mathematical content.

%=======================================================================
\appendix
\section{Theorem Index}\label{app:index}
%=======================================================================

The following table maps paper theorems to their Lean names and source files.
All results are fully machine-checked with zero \lean{sorry}.

\begin{longtable}{@{}>{\raggedright\arraybackslash}p{0.14\textwidth}
  >{\raggedright\arraybackslash}p{0.40\textwidth}
  >{\raggedright\arraybackslash}p{0.38\textwidth}@{}}
\toprule
\textbf{Ref.} & \textbf{Lean name} & \textbf{Source file} \\
\midrule
\endfirsthead
\multicolumn{3}{c}{\textit{(continued)}} \\
\toprule
\textbf{Ref.} & \textbf{Lean name} & \textbf{Source file} \\
\midrule
\endhead
\midrule
\multicolumn{3}{r}{\textit{continued on next page}} \\
\endfoot
\bottomrule
\endlastfoot

Def.~\ref{def:fiber} &
  \lean{ExtensionFiber} &
  \lean{FiberArchitecture.lean} \\

Thm.~\ref{thm:fiber-nonempty} &
  \lean{extension\_fiber\_nonempty} &
  \lean{FiberArchitecture.lean} \\

Thm.~\ref{thm:hasri} &
  \lean{rightHom\_surjective\_hasRightInverse} &
  \lean{FiberArchitecture.lean} \\

Def.~\ref{def:kernel-act} &
  \lean{kernelActOnFiber} &
  \lean{FiberArchitecture.lean} \\

Thm.~\ref{thm:free} &
  \lean{kernelAct\_free} &
  \lean{FiberArchitecture.lean} \\

Thm.~\ref{thm:transitive} &
  \lean{kernelAct\_transitive} &
  \lean{FiberArchitecture.lean} \\

Thm.~\ref{thm:section-exists} &
  \lean{section\_exists} &
  \lean{FiberArchitecture.lean} \\

Thm.~\ref{thm:split-char} &
  \lean{splitting\_iff\_section\_is\_homomorphism} &
  \lean{FiberArchitecture.lean} \\

Def.~\ref{def:cocycle} &
  \lean{sectionCocycle} &
  \lean{SchurZassenhaus.lean} \\

Thm.~\ref{thm:cocycle-spec} &
  \lean{sectionCocycle\_spec} &
  \lean{SchurZassenhaus.lean} \\

Thm.~\ref{thm:defect} &
  \lean{section\_mul\_eq} &
  \lean{SchurZassenhaus.lean} \\

Thm.~\ref{thm:split-criterion} &
  \lean{splits\_iff\_trivial\_cocycle} &
  \lean{SchurZassenhaus.lean} \\

Def.~\ref{def:section-diff} &
  \lean{sectionDiff} &
  \lean{CocycleCoboundary.lean} \\

Thm.~\ref{thm:section-relation} &
  \lean{section\_eq\_inl\_diff\_mul} &
  \lean{CocycleCoboundary.lean} \\

Thm.~\ref{thm:independence} &
  \lean{splitting\_independent\_of\_section} &
  \lean{CocycleCoboundary.lean} \\

Thm.~\ref{thm:any-trivial} &
  \lean{splits\_of\_any\_trivial\_cocycle} &
  \lean{CocycleCoboundary.lean} \\

Def.~\ref{def:conj-act} &
  \lean{sectionConjAct} &
  \lean{ConjugationAction.lean} \\

Thm.~\ref{thm:conj-spec} &
  \lean{inl\_sectionConjAct} &
  \lean{ConjugationAction.lean} \\

Thm.~\ref{thm:conj-inl} &
  \lean{conjAct\_inl} &
  \lean{ConjugationAction.lean} \\

Thm.~\ref{thm:conj-independent} &
  \lean{sectionConjAct\_independent} &
  \lean{ConjugationAction.lean} \\

Thm.~\ref{thm:cocycle-identity} &
  \lean{sectionCocycle\_isMulCocycle\textsubscript{2}\_conj} &
  \lean{CocycleIdentity.lean} \\

Thm.~\ref{thm:cocycle-splitting} &
  \lean{sectionCocycle\_of\_splitting} &
  \lean{CocycleIdentity.lean} \\

Def.~\ref{def:direct-product} &
  \lean{directProductExtension} &
  \lean{LocalGlobal.lean} \\

Thm.~\ref{thm:direct-splits} &
  \lean{directProduct\_splits}, \lean{directProduct\_has\_trivial\_cocycle} &
  \lean{LocalGlobal.lean} \\

Thm.~\ref{thm:obstruction} &
  \lean{not\_splits\_of\_all\_cocycles\_nontrivial} &
  \lean{LocalGlobal.lean} \\

Def.~\ref{def:pullback} &
  \lean{pullbackSubgroup} &
  \lean{LocalGlobal.lean} \\

Def.~\ref{def:restrict} &
  \lean{restrictExtension} &
  \lean{LocalGlobal.lean} \\

Thm.~\ref{thm:global-restricts} &
  \lean{splitting\_restricts}, \lean{global\_split\_implies\_local\_split} &
  \lean{LocalGlobal.lean} \\

Thm.~\ref{thm:local-obstruction} &
  \lean{local\_nonsplit\_obstructs\_global} &
  \lean{LocalGlobal.lean} \\

Thm.~\ref{thm:descent} &
  \lean{faithfullyFlat\_reflects\_injective}, etc. &
  \lean{FaithfullyFlatDescent.lean} \\

\end{longtable}

\noindent
All source files are located under
\lean{InfinityCompression/GeneralMethod/GroupExtension/} (group extension modules)
and \lean{InfinityCompression/GeneralMethod/Descent/} (faithfully flat descent).

%=======================================================================
\section{Module Dependency Graph}\label{app:deps}
%=======================================================================

The six core group extension modules (presented in this paper) form a layered
dependency structure:

\begin{center}
\begin{tabular}{@{}l@{\;\;$\longrightarrow$\;\;}l@{}}
\lean{FiberArchitecture.lean} & \lean{SchurZassenhaus.lean} \\
\lean{FiberArchitecture.lean} & \lean{ConjugationAction.lean} \\
\lean{SchurZassenhaus.lean} & \lean{CocycleCoboundary.lean} \\
\lean{SchurZassenhaus.lean} & \lean{CocycleIdentity.lean} \\
\lean{ConjugationAction.lean} & \lean{CocycleIdentity.lean} \\
\lean{FiberArchitecture.lean} & \lean{LocalGlobal.lean} \\
\lean{SchurZassenhaus.lean} & \lean{LocalGlobal.lean} \\
\end{tabular}
\end{center}

\noindent
\lean{FiberArchitecture.lean} is the root; it depends only on
\lean{Mathlib.GroupTheory.GroupExtension.Basic} and
\lean{Mathlib.Logic.Function.Basic}.  Three additional modules not discussed in
this paper (\lean{NonAbelianObstruction.lean}, \lean{SchurMultiplier.lean},
\lean{MathlibAdapter.lean}) explore further directions and adapter patterns.
The descent module \lean{FaithfullyFlatDescent.lean} is independent, depending
only on \lean{Mathlib.RingTheory.Flat.FaithfullyFlat.Descent}.

%=======================================================================
\begin{thebibliography}{99}
%=======================================================================

\bibitem{Lean4}
L.~de Moura and S.~Ullrich.
\newblock The {L}ean~4 theorem prover and programming language.
\newblock In \emph{Automated Deduction --- CADE 28}, volume 12699 of
  \emph{Lecture Notes in Computer Science}, pages 625--635. Springer, 2021.

\bibitem{Mathlib}
The Mathlib community.
\newblock The {L}ean mathematical library.
\newblock In \emph{Proceedings of the 9th ACM SIGPLAN International Conference
  on Certified Programs and Proofs (CPP)}, pages 367--381, 2020.
\newblock \textsc{doi}: \href{https://doi.org/10.1145/3372885.3373824}{10.1145/3372885.3373824}.
\newblock Mathlib~4 repository: \url{https://github.com/leanprover-community/mathlib4}.

\bibitem{Yamazaki2024}
Y.~Yamazaki.
\newblock Group extensions in {M}athlib.
\newblock Contribution to Mathlib~4, 2024.
\newblock \url{https://leanprover-community.github.io/mathlib4_docs/Mathlib/GroupTheory/GroupExtension/Defs.html}.

\bibitem{Kull2024}
S.~Kull.
\newblock Low-degree group cohomology in {M}athlib.
\newblock Contribution to Mathlib~4, 2024.
\newblock \url{https://leanprover-community.github.io/mathlib4_docs/Mathlib/RepresentationTheory/GroupCohomology/LowDegree.html}.

\bibitem{Livingston2024}
A.~Livingston.
\newblock Group cohomology in {M}athlib.
\newblock Contribution to Mathlib~4, 2022--2024.
\newblock \url{https://leanprover-community.github.io/mathlib4_docs/Mathlib/RepresentationTheory/GroupCohomology/Basic.html}.

\bibitem{Gonthier2013}
G.~Gonthier, A.~Asperti, J.~Avigad, Y.~Bertot, C.~Cohen, F.~Garillot,
  S.~Le~Roux, A.~Mahboubi, R.~O'Connor, S.~Ould~Biha, I.~Pasca, L.~Rideau,
  A.~Solovyev, E.~Tassi, and L.~Th\'ery.
\newblock A machine-checked proof of the odd order theorem.
\newblock In \emph{Interactive Theorem Proving (ITP 2013)}, volume 7998 of
  \emph{Lecture Notes in Computer Science}, pages 163--179. Springer, 2013.

\bibitem{HoTTBook}
The Univalent Foundations Program.
\newblock \emph{Homotopy Type Theory: Univalent Foundations of Mathematics}.
\newblock Institute for Advanced Study, 2013.
\newblock \url{https://homotopytypetheory.org/book/}.

\bibitem{Brown1982}
K.~S.~Brown.
\newblock \emph{Cohomology of Groups}, volume~87 of \emph{Graduate Texts in
  Mathematics}.
\newblock Springer, 1982.

\bibitem{Weibel1994}
C.~A.~Weibel.
\newblock \emph{An Introduction to Homological Algebra}, volume~38 of
  \emph{Cambridge Studies in Advanced Mathematics}.
\newblock Cambridge University Press, 1994.

\bibitem{Robinson1996}
D.~J.~S.~Robinson.
\newblock \emph{A Course in the Theory of Groups}, volume~80 of \emph{Graduate
  Texts in Mathematics}.
\newblock Springer, 2nd edition, 1996.

\bibitem{Rotman2009}
J.~J.~Rotman.
\newblock \emph{An Introduction to Homological Algebra}.
\newblock Springer, 2nd edition, 2009.

\end{thebibliography}

\end{document}
